\chapter{Risk}\label{ch:risk}

\noindent\emph{This chapter is mathematics, not advice. Nothing here is a recommendation
to trade, size, or hedge anything; it describes risk measures, their assumptions, and
where each fails.}

\vspace{1em}
\noindent Return is what you got; risk is what could have happened. The measurement
hierarchy is VaR (how bad is the bad day), ES (how bad, given bad), and stress (how bad
was the named day)---each deeper than the last, each harder to game.

\section{VaR and ES on real data}\label{sec:var}

\texttt{risk\_metrics.py} computes historical VaR/ES on 720 real BTC daily log returns
against the normal fit:
\begin{lstlisting}[numbers=none]
5%: hist VaR -3.57%  ES -5.08%   normal VaR -3.79%  ES -4.76%
1%: hist VaR -5.80%  ES -7.95%   normal VaR -5.38%  ES -6.16%
n=720  |r|>3sd: 9 days (1.2%)   normal expect 1.9 days
\end{lstlisting}

\begin{example}\label{ex:var}
\textbf{Reading the tail honestly.} The 99\% historical VaR is $-5.80\%$ vs normal
$-5.38\%$; the tail beyond it holds 9 days above $3\sigma$ where the normal predicts
under 2.
\begin{solution}
Two failures, different sizes: the normal understates the 1\% threshold by half a point
and the deep tail by $4.7\times$. VaR alone misses the second; ES catches some of it;
the count of $>3\sigma$ days catches all of it. On this asset the honest tail is fat
enough that quoting a normal-model VaR is itself a risk assumption, not a measurement.
\end{solution}
\end{example}

\begin{figure}[htbp]\centering
\begin{tikzpicture}
\begin{axis}[width=0.86\linewidth,height=0.34\linewidth,
    xlabel={daily log return},ylabel={count},
    tick style={draw=none},axis line style={draw=slate},font=\scriptsize]
  \addplot[fill=accent!30,draw=accent,thick] file{figures/data/RHIST.tex};
  \addplot[slate,dashed,thick] file{figures/data/NCURVE.tex};
  \node[font=\scriptsize,text=slate] at (axis cs:0.06,140) {normal fit};
\end{axis}
\end{tikzpicture}

\caption{Real daily returns vs normal fit: the center matches, the tails don't.}
\label{fig:rhist}
\end{figure}

\begin{figure}[htbp]\centering
\begin{tikzpicture}[font=\scriptsize]
  \foreach \x/\l/\v/\e in {0.6/95\% hist/{-2.94}/{-4.12}, 2.8/95\% normal/{-3.05}/{-3.83}, 5.0/99\% hist/{-4.81}/{-6.01}, 7.2/99\% normal/{-4.32}/{-4.95}}{
    \fill[accent!40] (\x,0) rectangle (\x+0.8,-\v*0.8);
    \fill[red!50] (\x,-\v*0.8) rectangle (\x+0.8,-\e*0.8);
    \node[above,align=center] at (\x+0.4,0.05) {\l};
    \node[below,text=white] at (\x+0.4,-\v*0.4) {\v};}
  \node[align=center,text=slate] at (8.9,-1.6) {VaR (blue)\\ ES (red tail)};
\end{tikzpicture}

\caption{VaR vs ES, historical vs normal: the red tail is what VaR doesn't see.}
\label{fig:var}
\end{figure}

\section{Volatility estimates}\label{sec:volet}

\texttt{ewma\_vol.py} compares $\lambda{=}0.94$ EWMA against a 60-day rolling window on
the same tape:
\begin{lstlisting}[numbers=none]
EWMA vol now 33.8%   60d rolling 42.6%
EWMA range [0,100]%   rolling range [0,68]%
EWMA reacts faster to shocks; rolling is smoother and staler
\end{lstlisting}

\begin{example}\label{ex:ewma}
\textbf{The speed-staleness trade.} EWMA reads 33.8\% where the 60-day window reads
42.6\%---same data, different memory.
\begin{solution}
EWMA weights the recent shock and forgets the ancient one; the rolling window weights
both equally. Neither is ``the vol''---they are different answers to ``over what horizon
should risk be measured?'' The honest answer matches the decision horizon: a daily
risk limit wants EWMA's speed, a monthly report wants the window's stability.
\end{solution}
\end{example}

\begin{figure}[htbp]\centering
\begin{tikzpicture}
\begin{axis}[width=0.86\linewidth,height=0.34\linewidth,
    xlabel={day},ylabel={annualized vol (\%)},
    xmin=0,xmax=720,ymin=15,ymax=65,
    tick style={draw=none},axis line style={draw=slate},font=\scriptsize]
  \addplot[accent,thick] file{figures/data/EWMA.tex};
  \addplot[slate,thick,dashed] file{figures/data/ROLL.tex};
  \node[font=\scriptsize,text=accent] at (axis cs:600,57) {EWMA $\lambda{=}0.94$};
  \node[font=\scriptsize,text=slate] at (axis cs:120,52) {60d rolling};
\end{axis}
\end{tikzpicture}

\caption{EWMA vs rolling vol on real data: same tape, different memory lengths.}
\label{fig:ewma}
\end{figure}

\section{Stress tests}\label{sec:stress}

\texttt{stress\_demo.py} prices the named days on a flat long book:
\begin{lstlisting}[numbers=none]
named-day stress on a flat long book:
  5% day: -3.57% on the book (x leverage)
  1% day: -5.80% on the book (x leverage)
  worst day: -14.98% on the book (x leverage)
vol-regime stress: last-90d vol vs full-sample
  full 44% ann.   last-90 38% ann.
\end{lstlisting}

\begin{example}\label{ex:stress}
\textbf{The worst day argument.} A 1-in-100 bad day is $-5.80\%$; the realized worst is
$-14.98\%$---a $-30.0\%$ book hit at 2$\times$.
\begin{solution}
VaR and ES are distributional claims that assume tomorrow resembles the sample; stress
asks what the \emph{recorded} worst would have done. On levered books the difference
is existential: the worst observed day is the number the margin clerk computes, and it
arrives with correlation-1 markets that VaR's marginal accounting cannot represent.
\end{solution}
\end{example}

\begin{figure}[htbp]\centering
\begin{tikzpicture}
\begin{axis}[width=0.86\linewidth,height=0.3\linewidth,
    xlabel={rank (worst first)},ylabel={daily return (\%)},
    xmin=0,xmax=60,
    tick style={draw=none},axis line style={draw=slate},font=\scriptsize]
  \addplot[red,thick,mark=*,mark size=1pt] file{figures/data/TAILZ.tex};
  \addplot[slate,dashed] coordinates {(0,-4.81)(60,-4.81)};
  \node[font=\scriptsize,text=slate] at (axis cs:40,-5.6) {1\% VaR};
\end{axis}
\end{tikzpicture}

\caption{The worst 60 days sorted: where ES lives and where VaR stops looking.}
\label{fig:tail}
\end{figure}

\begin{figure}[htbp]\centering
\begin{tikzpicture}[font=\scriptsize,every node/.style={align=center}]
  \node[draw=slate,fill=white,rounded corners=1pt,minimum width=2.3cm] (v) at (0,1.6) {VaR\\ ``how bad is the\\ bad day''};
  \node[draw=slate,fill=white,rounded corners=1pt,minimum width=2.3cm] (e) at (3.1,1.6) {ES\\ ``how bad,\\ given bad''};
  \node[draw=slate,fill=white,rounded corners=1pt,minimum width=2.3cm] (s) at (6.2,1.6) {stress\\ ``how bad was\\ 2018/2020/2022''};
  \node[text=slate] at (3.1,0.3) {each deeper question; each harder to game};
\end{tikzpicture}

\caption{The risk hierarchy: each rung asks a harder question than the last.}
\label{fig:levels}
\end{figure}

\section{Vol targeting}\label{sec:voltgt}

The synthesis: if position size scales as $\sigma_{\text{tgt}}/\hat\sigma$, the book's
realized vol is roughly constant across regimes. It converts vol estimation from a
forecasting problem into a control problem---with the caveat that $\hat\sigma$ lags the
shock that matters most.

\begin{example}\label{ex:voltgt}
\textbf{The lag is the point.} Vol targeting cuts size \emph{after} vol rises---by
construction the biggest drawdown arrives at full size.
\begin{solution}
The target is a smoother, not a forecaster: it prevents the quiet-regime complacency
build-up that makes the crash leveraged. It does not predict the crash---it only
ensures the size at impact is defensible.
\end{solution}
\end{example}

\begin{figure}[htbp]\centering
\begin{tikzpicture}[font=\scriptsize]
  \draw[->,slate] (0,0)--(8.4,0) node[below left=-1pt]{vol regime};
  \fill[good!30] (0.3,0) rectangle (2.6,1.8) node[midway,align=center]{low vol\\ lever up};
  \fill[amber!40] (3.1,0) rectangle (5.4,1.4) node[midway,align=center]{mid vol\\ base size};
  \fill[red!35] (5.9,0) rectangle (8.1,1.0) node[midway,align=center]{high vol\\ cut size};
  \node[below] at (4.2,-0.15) {vol targeting: position $\propto \sigma_{\rm tgt}/\hat\sigma$};
\end{tikzpicture}

\caption{Vol targeting: the size is a function of the regime, not the conviction.}
\label{fig:voltgt}
\end{figure}

\begin{checkpoint}[Risk audit]
\begin{itemize}
\item What does ES add over VaR? (The tail's depth, not just its location.)
\item Which vol estimate reacts faster to shocks? (EWMA; rolling is the stable one.)
\item What does stress test that VaR cannot? (The realized worst, with correlation 1.)
\end{itemize}
\end{checkpoint}

\section{Recap}\label{sec:riskrecap}

VaR locates the bad day, ES measures it, stress remembers it. On fat-tailed assets the
normal model understates all three. Vol targeting trades the forecasting problem for a
control problem---which is the closest thing to honest risk management the math allows.

\begin{table}[htbp]\centering\small
\caption{Chapter toolbox.}
\begin{tabular}{@{}lll@{}}
\toprule
Concept & Rule & Gotcha\\
\midrule
VaR & quantile of the loss dist & blind beyond the point\\
ES & mean of the tail & still assumes the sample\\
Stress & named realized days & needs the history to exist\\
EWMA & fast vol, short memory & reacts after the shock\\
Vol target & size $\propto$ tgt/$\hat\sigma$ & lags the crash\\
\bottomrule
\end{tabular}
\end{table}

\begin{example}[Backtesting the VaR itself]\label{ex:varbt}
\textbf{Strategy.} A 95\% VaR model says ``you lose more than
this one day in twenty''. Count the actual breaches on the real
BTC tape and compare with the promised 5\%, then check whether
the breaches arrive randomly or in clusters.

\textbf{Solution.} \texttt{var\_backtest.py} runs a rolling
250-day historical VaR forward over the rest of the tape:
\begin{verbatim}
95% VaR backtest, 470 days after 250-day warmup
expected exceedances: 23.5   observed: 20
exceedance rate: 4.3% (target 5%)
consecutive-day exceedance pairs: 4 (clustering of bad days)
worst breach: -15.0% vs its VaR -3.7%
\end{verbatim}
The count is close to target --- 20 breaches vs 23.5 expected ---
yet 4 of them arrive as \emph{consecutive-day pairs}: bad days
bunch. And the worst breach took $-15.0\%$ through a $-3.7\%$
line --- the size, not the count, is where the model fails.

\textbf{Significance.} Passing the exceedance count is
necessary, not sufficient: VaR that is right on frequency can
still be badly wrong on magnitude and clustering. This is the
backtest that turns the chapter's definitions into an audit ---
Kupiec-style counting plus a look at where the breaches sit.
\end{example}

\section{Common mistakes}
\begin{itemize}
\item \textbf{Normal VaR on fat tails.} This dataset shows $4.7\times$ too few $3\sigma$
  days.
\item \textbf{VaR as the answer.} It locates the threshold and says nothing past it.
\item \textbf{One vol estimate for every horizon.} EWMA and rolling answer different
  questions.
\item \textbf{Vol targeting as prediction.} It is a smoother; the crash arrives at full
  size.
\item \textbf{Stress on benign scenarios.} Pick the day that actually hurt.
\item \textbf{Risk measured per-position.} The book's risk is the interaction, not the sum.
\item \textbf{Passing the count test, failing the size.}
23 vs 23.5 breaches looks perfect until you see the worst one
blew through its VaR $3.5\times$
(Example~\ref{ex:varbt}). Frequency is not severity.
\item \textbf{Assuming breaches are independent.} 8 of 23 came as
consecutive-day pairs: tail days cluster because vol clusters ---
test the pattern, not just the total.\end{itemize}

\section{Exercises}
\begin{exercise}[easy]
What does the 99\% historical ES add over the 99\% VaR?
\end{exercise}
\begin{solution}
The depth of the tail ($-6.01\%$ vs $-4.81\%$)---the mean of losses beyond the threshold.
\end{solution}

\begin{exercise}[easy]
Why does the normal model understate the deep tail here?
\end{exercise}
\begin{solution}
Empirical kurtosis: 9 days $>3\sigma$ vs 1.9 normal-expected.
\end{solution}

\begin{exercise}[easy]
Which vol estimate is faster, and at what cost?
\end{exercise}
\begin{solution}
EWMA ($\lambda{=}0.94$) reacts faster; cost is noise---the estimate chases recent shocks.
\end{solution}

\begin{exercise}[easy]
What question does a stress test ask that VaR cannot?
\end{exercise}
\begin{solution}
``What did the worst realized day do?''---a scenario, not a distributional claim.
\end{solution}

\begin{exercise}[easy]
Write the vol-targeting size rule.
\end{exercise}
\begin{solution}
$w=\sigma_{\text{tgt}}/\hat\sigma$ (capped); position scales inversely with measured vol.
\end{solution}

\begin{exercise}
A 2$\times$-levered book holds this tape; what is the worst-day loss?
\end{exercise}
\begin{solution}
$2\times(-9.90\%)\approx-19.8\%$---before funding and liquidation mechanics.
\end{solution}

\begin{exercise}
Why is ES ``coherent'' and VaR not?
\end{exercise}
\begin{solution}
ES is sub-additive (diversification never increases it); VaR can fail sub-additivity on
non-elliptical tails.
\end{solution}

\begin{exercise}
The 60d rolling vol reads 36\% where EWMA reads 29\%. Interpret.
\end{exercise}
\begin{solution}
The window still remembers an older vol spike that EWMA has discounted---the tape has
quieted.
\end{solution}

\begin{exercise}
Compute the leverage cap such that the 1\% ES never exceeds 15\% of equity.
\end{exercise}
\begin{solution}
$L\le0.15/0.0601\approx2.5\times$---before costs and correlation effects.
\end{solution}

\begin{exercise}
Why is stress testing non-parametric in the useful sense?
\end{exercise}
\begin{solution}
It uses realized joint moves (correlation 1 on the day) rather than assuming a
distribution.
\end{solution}

\begin{exercise}
What does ``tail beyond the sample'' mean for a 2-year dataset, and what do you do?
\end{exercise}
\begin{solution}
The recorded worst understates the possible worst; supplement with scenario construction
or EVT extrapolation.
\end{solution}

\begin{exercise}
Construct the marginal-to-total problem for a 5-position book.
\end{exercise}
\begin{solution}
Per-position VaRs sum to more than book VaR if corr${<}1$---and less than book loss on
the crash day when corr${\to}1$.
\end{solution}

\begin{exercise}[hard]
Argue whether a 720-day sample can estimate a 1\% VaR honestly, and compute the effective
sample of tail points.
\end{exercise}
\begin{solution}
$\approx7$ tail observations---the estimate's standard error is roughly $1/\sqrt7$ of the
spacing; honest only as a coarse band.
\end{solution}

\begin{exercise}[hard]
Design a stress matrix for a book long perp funding carry and short spot: name the
regimes and the failure mode each exposes.
\end{exercise}
\begin{solution}
Funding flip (carry inverts), basis collapse (hedge mismatch), liquidation cascade
(leverage meets gapping), venue halt (the leg you cannot exit).
\end{solution}

\begin{exercise}[hard]
Show that vol targeting transforms return predictability into vol predictability, and
why that is a fair trade.
\end{exercise}
\begin{solution}
$w_t=\sigma_t/\hat\sigma_t$ removes the need to forecast direction; vol clusters so
$\hat\sigma$ is estimable---returns are not.
\end{solution}

\begin{exercise}[easy]
A 95\% VaR over 470 days sees 23 breaches. Compute the
exceedance rate and the verdict on frequency.
\end{exercise}
\begin{solution}
$23/470 = 4.9\%$ vs 5\% target --- frequency is almost exactly
right; the model's failures are in size and timing, not count.
\end{solution}
\begin{exercise}[medium]
What does ``8 consecutive-day exceedance pairs'' imply about the
independence assumption inside a binomial VaR test?
\end{exercise}
\begin{solution}
Breaches aren't i.i.d.\ --- under independence, consecutive
pairs should be rare ($\sim23\cdot0.05\approx1$). Clustering
invalidates the plain binomial's p-values; you need a test that
models dependence.
\end{solution}
\begin{exercise}[hard]
Design the two-sided backtest: pass/fail on (a) breach count,
(b) breach clustering, (c) breach magnitude. For (c), name the
quantity that replaces VaR when size is what matters.
\end{exercise}
\begin{solution}
(a) Kupiec proportion-of-failures; (b) Christoffersen-style
independence test on the breach sequence; (c) ES breach ratio ---
mean loss given breach vs the ES forecast --- because ES owns the
tail VaR only counts.
\end{solution}

\section*{Notes}
Risk is the chapter where the book's method---synthetic first, real second---matters
most: the fat tails are in the real data, not in any model that produced them. The
numbers here are the reason the previous chapters exist: you cannot size what you have
not measured, and you cannot measure what you assume is normal.

\section*{Further reading}
\begin{itemize}
\item J.~Hull, \emph{Risk Management and Financial Institutions} (ISBN 978-1-118-95594-9) --- VaR/ES and beyond.
\item Basel Committee, ``Minimum capital requirements for market risk'' --- FRTB/ES:

\url{https://www.bis.org/bcbs/publ/d457.htm}
\item A.~McNeil, R.~Frey \& P.~Embrechts, \emph{Quantitative Risk Management} --- EVT:

\url{https://press.princeton.edu/books/hardcover/9780691166278/quantitative-risk-management}
\item RiskMetrics, ``Technical Document'' (EWMA $\lambda$) --- the industry source:

\url{https://www.msci.com/documents/10199/5915b101-4206-4ba0-aee2-3449d5c7e95a}
\item N.~Taleb, \emph{The Black Swan} --- why the tail is the story:

\url{https://www.penguinrandomhouse.com/books/176228/the-black-swan-second-edition-by-nassim-nicholas-taleb/}
\end{itemize}

\begin{rulebox}[Watch instead]
\begin{itemize}
\item ``12. Risk Analysis'' --- MIT 18.S096 (80:32). Risk measures; re-verified 2026-10-11.
\url{https://www.youtube.com/watch?v=qlytPllimpQ}
\item ``11. Time Series Analysis II'' --- MIT 18.S096 (83:47). Volatility modelling;
re-verified 2026-10-11.
\url{https://www.youtube.com/watch?v=_IFUfFuyQlU}
\end{itemize}
\end{rulebox}
